% SOURCE_STATEMENT: "When both particle species have condensates, give a rigorous analysis of the relation between the two condensates."
% Preparation and final self-audit: CONTRIBUTING.md prompts C and D; source/audit artifacts retained in the local physics_statements_2026_10_11 batch.
% !TeX program = lualatex
\documentclass[11pt]{article}
\usepackage[margin=25mm]{geometry}
\usepackage{fontspec}
\setmainfont{Latin Modern Roman}
\usepackage{amsmath,amssymb,mathtools}
\usepackage{unicode-math}
\setmathfont{Latin Modern Math}
\usepackage{xurl,hyperref}
\hypersetup{colorlinks=true,urlcolor=blue}
\setcounter{secnumdepth}{0}
\setlength{\parindent}{0pt}
\setlength{\parskip}{5pt}
\setlength{\emergencystretch}{3em}
\newcommand{\sref}[2]{\hyperlink{source:#1:#2}{[S#2]}}
\newcommand{\cataloguescope}[1]{\paragraph{Recorded target.}#1}
\newcommand{\R}{\mathbb R}
\newcommand{\C}{\mathbb C}
\newcommand{\N}{\mathbb N}
\newcommand{\Z}{\mathbb Z}
\newcommand{\Q}{\mathbb Q}
\newcommand{\F}{\mathbb F}
\newcommand{\eps}{\varepsilon}
\begin{document}
% UnsolvedMath ID 30000226; source catalogue code OWR-827-001
\setcounter{section}{30000225}
\section{Relationship Between Two-Species Condensates}\label{30000226}
{\small UnsolvedMath ID 30000226 · Mathematical Physics}
\subsection{Definitions and mathematical statement}

\paragraph{Shared notation.}$\mathbb N=\{1,2,\ldots\}$; the standard number fields have their usual meanings. All additional conventions are stated locally.


\paragraph{Definitions and conventions.} The source concerns a two-species zero-range particle process, not Bose condensation. On $L$ sites, $\eta_i(x)\in\mathbb N_0$ is the occupation of species $i=1,2$ and $\sum_x\eta_i(x)=N_i$ is conserved. A zero-range jump rate depends on the occupations at the departure site. In the thermodynamic limit $N_i/L\to\rho_i$, a condensate is an occupation of order $L$ on a set of sites of vanishing volume fraction, containing the excess mass above the critical background density. The report's stationary canonical ensemble and weight assumptions determine which density pairs have one or two condensed species.

\paragraph{Source scope and qualifications.} The source proves ensemble and single-site condensation facts under assumptions on its weights. The remaining relation of two condensates includes their locations and coupled mass fluctuations; the brief source does not prescribe one particular statistic or a universal rate function.

\paragraph{Catalogue mathematical target.}
\paragraph{Catalogue wording retained for traceability.}
When both particle species have condensates, give a rigorous analysis of the relation between the two condensates.

\paragraph{Saved literature context.} The supplied catalogue's assessment dated 2026-08-21 records: Equivalence of ensembles and single-site condensation are proved under substantial assumptions, but the relation between two condensates of two species is still identified as an open rigorous problem. The assessment is retained as dated provenance; it is not a new comprehensive literature or open-status certification.
\subsection{Short English statement}
When both species condense in a two-species zero-range process, determine rigorously how their condensates are related.
\subsection{Sources}
\begin{itemize}
\item[\hypertarget{source:30000226:1}{S1}] ULAM.AI and the UnsolvedMath Contributors, UnsolvedMath, record 30000226 (OWR-827-001). Catalogue curation under CC BY 4.0; definitions and English exposition added locally.
\url{https://huggingface.co/datasets/ulamai/UnsolvedMath}.
\item[\hypertarget{source:30000226:2}{S2}] Mini-Workshop: Particle Systems with Several Conservation Laws: Fluctuations and Hydrodynamic Limit (2005), Oberwolfach Reports, DOI 10.4171/owr/2005/22, pp. 1199--1230; official report source and mathematical context.
\url{https://doi.org/10.4171/owr/2005/22}.
\item[\hypertarget{source:30000226:3}{S3}] Official EMS Press PDF of the same Oberwolfach report; PDF page 21, printed page 1219.
\url{https://ems.press/content/serial-article-files/45996}.
\item[\hypertarget{source:30000226:4}{S4}] Grosskinsky, two-species zero-range invariant measures; reference in the saved dated catalogue literature review.
\url{https://doi.org/10.1016/j.spa.2007.09.006}.
\end{itemize}
\label{end:30000226}
\end{document}
